% ===== latex-recomposed:document =====
% flm: style file flmrecomposed.7714d560.sty
\flmHeading[toc=true]
\section{Links in a title: \flmHref{https://example.com/a}{the site} and \flmUrl{https://example.com/b}}{\flmLabels{}}

A paragraph holds \flmHref{https://example.com}{a hyperlink}, an address \flmUrl{https://example.com/100\%25/a_b~c?q=1&r=2\#top} and \flmEmail{someone@example.com}; \flmTerm{ref1}{entropy} is defined below.

\begin{defterm}{Markov chain}\phantomsection\flmLabelsWithDisplayText{Markov chain}{\flmLabels{\label{ref2}\label{ref3}\label{ref4}}}\ignorespaces
A random process; a \flmTermInDefterm{Markov chain} forgets its past, as all \flmTermInDefterm{Markov chains} do, and may produce \flmTerm{ref1}{entropy}.
\end{defterm}

\begin{defterm}{entropy}\phantomsection\flmLabelsWithDisplayText{entropy}{\flmLabels{\label{ref1}\label{ref5}}}\ignorespaces
\emph{A measure} of disorder, which a \flmTerm{ref2}{process} may produce.
\end{defterm}

\begin{defterm}{\(\alpha\)-decay}\phantomsection\flmLabelsWithDisplayText{\(\alpha\)-decay}{\flmLabels{\label{ref6}}}\ignorespaces
A decay that emits an \(\alpha\) particle.
\end{defterm}

See \flmTerm{ref1}{entropy}, \flmTerm{ref1}{the entropy}, \flmTerm{ref3}{Markov chains}, \flmTerm{ref6}{\(\alpha\)-decay}, \flmRef{ref4} and \flmRef{ref5}; \flmHref{https://example.com/g}{a page on \flmTerm{ref1}{entropy}}.

\flmHeading[toc=true]
\section{A heading about \flmTerm{ref1}{disorder}}{\flmLabels{}}

\begin{itemize}
\item An item holds \flmHref{https://example.org}{a link} and \flmTerm{ref1}{entropy}.
\item\begin{defterm}{list term}\phantomsection\flmLabelsWithDisplayText{list term}{\flmLabels{\label{ref7}}}\ignorespaces
A term defined in a list item, where \flmTermInDefterm{list term} is itself.
\end{defterm}
\end{itemize}

A note holds links.\flmXNfootnote{The note holds \flmUrl{https://example.net/} and \flmTerm{ref1}{entropy}.} Another defines a term.\flmXNfootnote{\begin{defterm}{note term}\phantomsection\flmLabelsWithDisplayText{note term}{\flmLabels{\label{ref8}}}\ignorespaces
A term defined in a footnote.
\end{defterm}} See \flmTerm{ref8}{note term} and \flmTerm{ref7}{list term}.

\begin{equation}
S = k \log W \text{ where \flmHref{https://example.com/f}{a link} and \flmTerm{ref1}{entropy} stand}
\end{equation}

Through definitions: \flmUDseeentropy{}; \flmHref{https://doi.org/10.1000/xyz-1}{10.1000/xyz-1}; \flmUrl{https://example.com/c/\#sec-intro}; \flmHref{https://doi.org/10.1000/a\%20b\#c}{DOI}; \flmTerm{ref2}{this term}.

\begin{defterm}{expanded term}\phantomsection\flmLabelsWithDisplayText{expanded term}{\flmLabels{\label{ref9}}}\ignorespaces
A term defined through a user definition.
\end{defterm}


% flm: unsupported element flm.userdefs.macro.kdt

See \flmTerm{ref9}{expanded term}.%%
% ===== latex-recomposed:style flmrecomposed.7714d560.sty =====
\NeedsTeXFormat{LaTeX2e}
\ProvidesPackage{flmrecomposed.7714d560}
% --- packages ---
% flm.latex.pkg.amsmath
\RequirePackage{amsmath}
% flm.latex.pkg.hyperref
\RequirePackage{hyperref}
% flm.latex.pkg.zref-clever
\RequirePackage{zref-clever}
% --- commands and environments the library defines for the document ---
% flm.latex.def.xnotes.footnote
% \flmXNfootnote[<labels>]{<note>}: a note of the kind `footnote' in the FLM document;
% by default it passes the note, its labels first, to \footnote.
% Redefine it to set these notes another way.
\ProvideDocumentCommand\flmXNfootnote{+O{} +m}{\footnote{#1#2}}
% flm.latex.def.href
% Links to an address outside the document, one hook per FLM command:
% \flmHref{<address>}{<text>}: <text>, linked to <address>, as \href does.
% \flmUrl{<address>}: the display of <address>, linked to <address> and
%   typeset the way \url typesets an address.
% \flmEmail{<address>}: the display of mailto:<address> -- the address --
%   linked to mailto:<address> and typeset the same way.
% The display is what the other outputs of FLM show of such a link: its target
% without one leading http://, https:// or mailto: and without the trailing
% run of the characters / # ? -- https://example.com/?# shows example.com.
% Every argument is an ordinary one: in <address> the characters # % \ are
% written \# \% \\. Redefine a hook to link or typeset differently.
%
% The display is derived by \flm@href@display{<address>}, written in expl3,
% LaTeX's programming layer, which is part of the format and needs no package;
% its names read only between \ExplSyntaxOn and \ExplSyntaxOff. It reads the
% address as a string of characters (\tl_to_str:n), drops the backslash of
% each \# \% \\, then the prefix, then the trailing run, one character at a
% time. It works by expansion alone and leaves the display and nothing else,
% so that it serves where a hook is typeset and where hyperref builds a PDF
% bookmark by expansion. The macros with @ in their names are internals, not a
% customization point: redefine a hook, not its parts.
\ExplSyntaxOn
\def\flm@href@display#1{\exp_args:Ne \flm@href@trim
  {\exp_args:Ne \flm@href@unprefix
    {\exp_args:Ne \flm@href@unescape {\tl_to_str:n {#1}}}}}
% Every character, a backslash dropped and the character after it kept.
\def\flm@href@unescape#1{\str_if_empty:nF {#1}
  {\str_if_eq:eeTF {\str_head:n {#1}} {\c_backslash_str}
    {\exp_args:Ne \flm@href@unescape@one {\str_tail:n {#1}}}
    {\flm@href@unescape@one {#1}}}}
\def\flm@href@unescape@one#1{\str_head:n {#1}
  \exp_args:Ne \flm@href@unescape {\str_tail:n {#1}}}
% The first of the three prefixes the string begins with, taken off.
\def\flm@href@unprefix#1{\flm@href@unprefix@try {#1} {http://}
  {\flm@href@unprefix@try {#1} {https://}
    {\flm@href@unprefix@try {#1} {mailto:} {#1}}}}
\def\flm@href@unprefix@try#1#2#3{\str_if_eq:eeTF
  {\str_range:nnn {#1} {1} {\str_count:n {#2}}} {#2}
  {\str_range:nnn {#1} {\str_count:n {#2} + 1} {-1}} {#3}}
% The last character taken off while it is one of / # ?.
\def\flm@href@trim#1{\exp_args:Ne \flm@href@trim@last {\str_item:nn {#1} {-1}} {#1}}
\def\flm@href@trim@last#1#2{\bool_lazy_any:nTF
  {{\str_if_eq_p:nn {#1} {/}} {\str_if_eq_p:ee {#1} {\c_hash_str}} {\str_if_eq_p:nn {#1} {?}}}
  {\exp_args:Ne \flm@href@trim {\str_range:nnn {#2} {1} {-2}}} {#2}}
% The display of <target>, typeset by \nolinkurl.
\def\flm@href@show#1{\exp_args:Ne \nolinkurl {\flm@href@display {#1}}}
\ExplSyntaxOff
\ProvideDocumentCommand\flmHref{m m}{\href{#1}{#2}}
\ProvideDocumentCommand\flmUrl{m}{\href{#1}{\flm@href@show{#1}}}
\ProvideDocumentCommand\flmEmail{m}{\href{mailto:#1}{\flm@href@show{mailto:#1}}}
% What such a link contributes to a PDF bookmark: its text or its display
% alone. The hook runs when the document begins, so that it also holds where
% the document loads hyperref itself, after this file, having named it in the
% render option skip_packages; it does nothing where hyperref is absent.
\AtBeginDocument{\ifdefined\pdfstringdefDisableCommands
  \pdfstringdefDisableCommands{%
    \def\flmHref#1#2{#2}%
    \def\flmUrl#1{\flm@href@display{#1}}%
    \def\flmEmail#1{\flm@href@display{mailto:#1}}%
  }\fi}
% flm.latex.def.defterm
% Defined terms. A term definition is written as the environment defterm,
% whose one argument is the term. This file defines the environment further
% below, in the part headed "definitions written from the configuration of
% FLM features"; the environment's opening text is the lead-in that the
% document was recomposed with. The body of a term definition opens with
% \phantomsection, so that a link to the term definition leads to it, and
% with \flmLabelsWithDisplayText, so that a reference to one of its labels
% prints the term.
% \flmTerm{<label>}{<shown>}: <shown>, a link to the definition of a term,
%   <label> being a label of that definition, as \flmHyperref takes it.
% \flmTermInDefterm{<shown>}: <shown> written inside the definition of its
%   own term: the term being defined, set apart, with no link.
% Redefine either command to set term links another way.
\ProvideDocumentCommand\flmTerm{m m}{\flmHyperref{#1}{#2}}
\ProvideDocumentCommand\flmTermInDefterm{m}{\emph{#1}}
% What a term link contributes to a PDF bookmark: its text alone. The code
% below runs when the document begins (\AtBeginDocument) and does nothing
% where hyperref is not loaded.
\AtBeginDocument{\ifdefined\pdfstringdefDisableCommands
  \pdfstringdefDisableCommands{%
    \def\flmTerm#1#2{#2}%
    \def\flmTermInDefterm#1{#1}%
  }\fi}
% flm.latex.def.refs
% \flmRef[<options>]{<label>}: a reference that gives the name of its target
% ("Section 2"), set by \zcref of the zref-clever package.
%
% The options are the ones the document wrote, in FLM's own vocabulary. This
% definition translates each of them by its key -- a value written after the
% key is ignored, as in FLM -- and drops every key it has no translation for:
%
%   noprefix    -> noname                 the number or title alone, with no
%                                         name before it ("2", not
%                                         "Section 2")
%   S           -> S                      the first name capitalized and
%                                         unabbreviated, for a reference that
%                                         opens a sentence
%   capital     -> cap=true,abbrev=false  the names capitalized ("Lemma 1")
%   lowercase   -> cap=false,abbrev=false the names in lowercase ("lemma 1")
%   abbreviated -> abbrev=true            the abbreviated names ("Lem. 1")
%                                         where the reference type declares
%                                         them, the full names otherwise
%   anything else -> dropped              the default wording; a variant of
%                                         your own naming has no counterpart
%                                         in zref-clever, which would stop on
%                                         a key it does not know
%
% Redefine \flmRef to set references differently -- another package, another
% wording, your own variants, no link. Your definition sees the options as the
% document wrote them, in FLM's vocabulary, and maps them as it likes.
%
% The translation uses expl3, LaTeX's programming layer, which is part of the
% format and needs no package; its names read only between \ExplSyntaxOn and
% \ExplSyntaxOff. \keyval_parse:NNn reads the option list as LaTeX reads any
% key-value list: it hands a key written alone to the first macro below and a
% key with its value to the second, which drops the value. \str_case:nnF looks
% the key up in the table and appends what \zcref is to be given, or nothing.
% The macros with @ in their names are internals, not a customization point:
% redefine \flmRef, not its parts.
\ExplSyntaxOn
\def\flm@ref@add#1{\tl_put_right:Nn \flm@ref@options {#1,}}
\def\flm@ref@key#1{\str_case:nnF {#1}
  {{noprefix}{\flm@ref@add{noname}}
   {S}{\flm@ref@add{S}}
   {capital}{\flm@ref@add{cap=true,abbrev=false}}
   {lowercase}{\flm@ref@add{cap=false,abbrev=false}}
   {abbreviated}{\flm@ref@add{abbrev=true}}}
  {}}
\def\flm@ref@pair#1#2{\flm@ref@key{#1}}
\def\flm@ref@zcref#1#2{\zcref[#1]{#2}}
\ProvideDocumentCommand\flmRef{O{} m}{%
  \def\flm@ref@options{}%
  \keyval_parse:NNn \flm@ref@key \flm@ref@pair {#1}%
  \exp_args:No \flm@ref@zcref {\flm@ref@options} {#2}}
\ExplSyntaxOff
% flm.latex.def.hyperref
% \flmHyperref{<label>}{<text>}: <text>, shown as a link to the place <label>
% refers to.
%
% Making a link in a PDF takes the hyperref package, and loading it changes a
% whole document -- every reference becomes a link, bookmarks are built. That
% is the author's decision, not this file's, so this file neither loads
% hyperref nor asks for it: the definition below links through it where the
% document loaded it and prints the text alone where it did not, so the
% document compiles either way. The question is about this document's own
% preamble, so the answer is the same wherever the document is compiled; it is
% asked here, at the reference, so it is right whether this file was loaded
% before hyperref or after it.
%
% Redefine \flmHyperref to link differently -- another package, a footnote
% with the target's number, no link at all.
\ProvideDocumentCommand\flmHyperref{m m}{%
  \IfPackageLoadedTF{hyperref}{\hyperref[{#1}]{#2}}{#2}}
% What such a reference contributes to a PDF bookmark: the text alone. A
% bookmark is plain text and carries no link, and without this line the label
% would appear in it beside the text. The hook runs when the document begins
% and does nothing where hyperref is absent.
\AtBeginDocument{\ifdefined\pdfstringdefDisableCommands
  \pdfstringdefDisableCommands{%
    \def\flmHyperref#1#2{#2}%
  }\fi}
% flm.latex.def.label-display-text
% \flmLabelsWithDisplayText{<text>}{<labels>}: makes <labels> refer to this
% place, with <text> as what a reference to them prints. It is the only
% command for that: a heading the class did not number uses it, and so does
% anything else that has a name but no number.
%
% It places no anchor of its own. Under hyperref a sectioning command, starred
% or not, already places one at the title; a \label after it points there.
% \refstepcounter would place a second anchor here, below the title, and a
% link would lead below the heading. So the counter is stepped with
% the plain \stepcounter and the two things this needs are set here:
% \@currentlabel, which is what a \label records, and \@currentcounter,
% which is what zref-clever reads to give the label its reference type. The
% group keeps every change local.
%
% <text> is expanded to make \@currentlabel, and \label does not survive that
% expansion; so a \label inside <text> (for example a footnote's) is dropped
% from the label's text only. At its own place in the document it still makes
% its own label.
%
% A caller with no sectioning command before it -- for example a defined term
% in running text -- writes \phantomsection at its own place, so that a link
% has an anchor of its own to lead to.
%
% Redefine \flmLabelsWithDisplayText to change how such labels are made.
\ifcsname c@flm@internal@refcnt\endcsname\else\newcounter{flm@internal@refcnt}\fi
\ProvideDocumentCommand\flmLabelsWithDisplayText{m m}{%
  \begingroup
    \stepcounter{flm@internal@refcnt}%
    {\let\label\@gobble\protected@xdef\flm@lwdt@text{#1}}%
    \let\@currentlabel\flm@lwdt@text
    \def\@currentcounter{flm@internal@refcnt}%
    #2%
  \endgroup}
% flm.latex.def.zref-clever-type
% The zref-clever reference type of the counter \flmLabelsWithDisplayText
% steps, with every name prefix empty, so that a reference to such a label
% prints its text alone ("Introduction", not "Section Introduction").
% Every other label keeps the name its own type gives it. Redefine these
% settings, or the type of the counter, to put something before the text.
\zcsetup{countertype={flm@internal@refcnt=flm--internal--ref}}
\zcRefTypeSetup{flm--internal--ref}{Name-sg = , name-sg = , Name-pl = , name-pl = ,}
% flm.latex.def.labels
% \flmLabels{<\label commands>}: the labels of one place in the document.
% In text it runs them as written. On a line of a display-math
% environment, where amsmath keeps the last label of a line and loses the
% ones before it, it runs the first and records every further one as an alias:
% a line per alias is written to the .aux file after the last page and read on
% the next run, making the alias refer to what the first label refers to
% (for \ref and for \zcref alike). Redefine \flmLabels to change what a
% group of labels does; \def\flmLabels#1{#1} restores plain LaTeX behavior,
% amsmath's "Multiple \label's" message included.
\ProvideDocumentCommand\flmLabels{m}{%
  \ifmmode\expandafter\flm@labels@math\else\expandafter\@firstofone\fi{#1}}
% Reading the group: the first \label as written, the rest as aliases of it.
\def\flm@labels@math#1{\def\flm@lbl@first{}\flm@lbl@scan#1\flm@lbl@stop}
\def\flm@lbl@scan{\@ifnextchar\label\flm@lbl@take\flm@lbl@rest}
\def\flm@lbl@take\label#1{%
  \ifx\flm@lbl@first\@empty
    \def\flm@lbl@first{#1}%
    \label{#1}%
  \else
    \expandafter\flm@recordalias\expandafter{\flm@lbl@first}{#1}%
  \fi
  \flm@lbl@scan}
\def\flm@lbl@rest#1\flm@lbl@stop{#1}
% Collecting the aliases. A name is recorded once however often the line is
% typeset (amsmath measures a line before setting it).
\def\flm@aliaslist{}
\def\flm@recordalias#1#2{%
  \ifcsname flm@alias@#2\endcsname\else
    \expandafter\gdef\csname flm@alias@#2\endcsname{}%
    \g@addto@macro\flm@aliaslist{\flm@alias@entry{#2}{#1}}%
  \fi}
% Writing them, from the hook that runs after the last page has been shipped
% out: every \newlabel of that page has reached the .aux file by then, so the
% alias lines come after them. \AtEndDocument would run too early.
\def\flm@alias@write#1#2{\immediate\write\@auxout{\string\flm@aliaslabel{#1}{#2}}}
\AddToHook{enddocument/afterlastpage}{\let\flm@alias@entry\flm@alias@write\flm@aliaslist}
% Reading them back: the alias refers to what the original refers to. The
% assignments are global because the .aux file is read inside a group.
\def\flm@aliaslabel#1#2{%
  \ifcsname r@#2\endcsname
    \global\expandafter\let\csname r@#1\expandafter\endcsname\csname r@#2\endcsname
  \fi
  \ifcsname Z@R@#2\endcsname
    \global\expandafter\let\csname Z@R@#1\expandafter\endcsname\csname Z@R@#2\endcsname
  \fi}
% flm.latex.def.heading
% \flmHeading[<keys>]<sectioning command><star>{<title>}{<labels>}: a heading
% and the labels that refer to it, written on two lines for readability, as in
%   \flmHeading[toc=true]
%   \section*{Title}{\flmLabels{\label{ref1}}}
% It runs the sectioning command as written, then makes the labels refer to
% the heading: by its number when the class numbered it, by its title when it
% did not (a starred heading, or a level below secnumdepth). It finds out
% which without patching anything: a class numbers a heading through
% \refstepcounter, which records the counter's name in \@currentcounter.
%
% The optional argument is a key-value list. It is written here, on the hook,
% and never on the sectioning command: a starred sectioning command takes no
% optional argument, and giving one to \section* is not an error, it is
% silently wrong -- \section*[Short]{Long} takes one token as its heading,
% so the opening bracket alone becomes the heading and Short, the closing
% bracket and Long all run together into the following paragraph. Two keys
% are defined:
%
%   toc=true|false     whether the heading is listed in the table of contents.
%                      Absent, it is listed as LaTeX lists it. After a starred
%                      sectioning command with toc=true, which LaTeX never
%                      lists, \addcontentsline{toc}{<level>}{<entry>} adds it:
%                      <level> is the command's own name (section for
%                      \section), <entry> the alternative title if alt was
%                      written, else the title; under hyperref the starred
%                      command has already placed the anchor the line links
%                      to. On an unstarred heading toc=false is not honored:
%                      the heading stays listed and a warning says so.
%
%   alt={Short title}  an alternative title: what goes into the table of
%                      contents, the running head and the PDF bookmark, while
%                      the page keeps the full title. It is handed to the
%                      sectioning command as its own optional argument, so it
%                      does what \section[Short title]{Full title} does, and
%                      nothing else: a reference to an unnumbered heading still
%                      prints the full title. A starred heading does not hand
%                      it to the sectioning command; where it is written
%                      toc=true, alt is its line in the table of contents.
%
% What counts is whether the key was written, not what it holds. alt={} is an
% empty alternative title and is handed over as \section[]{Full title}, which
% is how LaTeX asks for an empty table-of-contents entry; no optional argument
% at all, and an empty one ([]), hand nothing over and leave the sectioning
% command exactly as it was. Writing alt with no value is LaTeX's own error,
% as is a key that is not defined here, and a value of toc other than true or
% false. The values and the record that alt was written are cleared at every
% call, so one heading's keys never reach the next. The recomposed LaTeX writes
% toc on every heading and alt where the heading was given an alternative
% title; both keys are also there for a document written by hand and for a
% user who redefines the hook.
%
% Redefine \flmHeading to change what a heading does. The block below uses
% expl3, LaTeX's programming layer, which is part of the format and needs no
% package; its names read only between \ExplSyntaxOn and \ExplSyntaxOff.
% \keys_define:nn declares the keys and \keys_set:nn reads a list of them;
% \bool_if:NTF asks whether the key was written. The three declarations are
% guarded on their own names, as every declaration in this file is, so that two
% of these style files can be loaded in one LaTeX run. The macros with @ in
% their names are internals, not a customization point: redefine \flmHeading,
% not its parts.
\ExplSyntaxOn
\ifcsname flm@heading@alt\endcsname\else\tl_new:N \flm@heading@alt\fi
\ifcsname flm@heading@given\endcsname\else\bool_new:N \flm@heading@given\fi
\ifcsname flm@heading@toc\endcsname\else\tl_new:N \flm@heading@toc\fi
\keys_define:nn {flm/heading}
  {alt .code:n = {\tl_set:Nn \flm@heading@alt {#1}
                  \bool_set_true:N \flm@heading@given} ,
   alt .value_required:n = true ,
   toc .choices:nn = {true, false} {\tl_set_eq:NN \flm@heading@toc \l_keys_choice_tl} ,
   toc .value_required:n = true}
\def\flm@heading@withalt#1#2#3{#2[#1]{#3}}
% The contents line of a starred heading: <level> then <entry>.
\def\flm@heading@addtoc#1#2{\addcontentsline{toc}{#1}{#2}}
% toc=false on an unstarred heading: the heading stays listed.
\def\flm@heading@notoc#1{\PackageWarning{flm}{toc=false~on~the~unstarred~
  \string#1~is~not~honored~by~the~default~\string\flmHeading;~the~heading~
  stays~in~the~table~of~contents}}
\ProvideDocumentCommand\flmHeading{O{} m s m m}{%
  \tl_clear:N \flm@heading@alt
  \bool_set_false:N \flm@heading@given
  \tl_clear:N \flm@heading@toc
  \keys_set:nn {flm/heading} {#1}%
  \def\@currentcounter{}%
  \IfBooleanTF{#3}
    {#2*{#4}%
     \str_if_eq:VnT \flm@heading@toc {true}
       {\edef\flm@heading@level{\cs_to_str:N #2}%
        \bool_if:NTF \flm@heading@given
          {\exp_args:NVV \flm@heading@addtoc \flm@heading@level \flm@heading@alt}
          {\exp_args:NV \flm@heading@addtoc \flm@heading@level {#4}}}}
    {\str_if_eq:VnT \flm@heading@toc {false} {\flm@heading@notoc{#2}}%
     \bool_if:NTF \flm@heading@given
       {\exp_args:NV \flm@heading@withalt \flm@heading@alt #2{#4}}
       {#2{#4}}}%
  \ifx\@currentcounter\@empty
    \flmLabelsWithDisplayText{#4}{#5}%
  \else
    #5%
  \fi}
\ExplSyntaxOff
% --- definitions written from the configuration of FLM features ---
% defterm (the defterm feature's, written from its configuration): defined unless the name is already taken, so a definition of it made before this file is loaded wins; \RenewDocumentEnvironment after loading this file changes it
\ProvideDocumentEnvironment{defterm}{+m}{}{}
% --- user definitions ---
% \flmUDseeentropy (the document's \seeentropy): defined unless the name is already taken, so a definition of it made before this file is loaded wins; \renewcommand after loading this file changes it
\providecommand\flmUDseeentropy{see \flmTerm{ref1}{entropy}}
\endinput
